% Приложения
\clearpage
\appendix
\renewcommand{\thesection}{\Asbuk{section}}
\setcounter{section}{0}
% В приложении все \subsubsection* — без номера.
% \titleformat* меняет только шрифт, но оставляет before-code (где лежит \thesubsubsection).
% Нужно переопределить формат целиком, убрав счётчик из before-code:
\titleformat{\subsubsection}[block]{\normalfont\large\bfseries\raggedright}{}{0em}{#1}

\section{Исходный код бенчмарка Loop/Select}
\label{app:bench-src}

\lstset{language=Refal}
\begin{lstlisting}[
    caption={Бенчмарк Loop/Select (bench/loop-select.ref)},
    label={lst:loop-select-src}]
* Loop/Select benchmark - demonstrates O(N) for bootstrapped deque
* Run: java -jar target/refal5q.jar bench/loop-select.ref 10000

$ENTRY Go {
  = <TimeElapsed 0>
    <Loop Left () <Numb <Arg 1>>>
    <Prout <TimeElapsed>>;
}

Loop {
  s.Side (e.Acc) 0 = /* done */;
  s.Side (e.Acc) s.Rest
    = <Loop <Select s.Side (e.Acc A) (e.Acc B)> <- s.Rest 1>>;
}

Select {
  Left  t.Left t.Right = Right t.Left;
  Right t.Left t.Right = Left  t.Right;
}
\end{lstlisting}

\section{Инструкция по воспроизведению измерений}
\label{app:repro}

\subsubsection*{Предварительные требования}

\begin{itemize}
  \item Java~25~LTS и Apache~Maven~3.9.x (должны быть доступны в PATH);
  \item PowerShell~5.1 или новее;
  \item для сравнения с Refal-5~PZ: дистрибутив Refal-5~PZ с исполняемыми
        \prog{refc.exe} и \prog{refgo.exe};
  \item для сравнения с Refal-5J: собранный \prog{r5jc.jar} и \prog{r5jrt.jar}.
\end{itemize}

\subsubsection*{Сборка интерпретатора}

\begin{longlisting}
\caption{Сборка и запуск бенчмарка}\label{lst:repro-build}
\begin{minted}{bash}
cd refal5j-impl
mvn package -DskipTests

# our implementation only:
powershell -File bench\loop-select-bench.ps1

# with Refal-5 PZ:
powershell -File bench\loop-select-bench.ps1 -WithPZ

# with all three:
powershell -File bench\loop-select-bench.ps1 -WithPZ -WithR5J
\end{minted}
\end{longlisting}

Результаты сохраняются в \prog{bench/loop-select-results.csv}.

\subsubsection*{Построение графика}

\begin{longlisting}
\caption{Построение графика}\label{lst:repro-plot}
\begin{minted}{bash}
pip install matplotlib numpy
python bench\plot-benchmark.py --out bench\loop-select-graph.pdf
\end{minted}
\end{longlisting}

\subsubsection*{Настройка путей}

Пути к исполняемым файлам Refal-5~PZ и Refal-5J задаются переменными
\prog{\$Refal5}, \prog{\$R5JBin} и \prog{\$R5JLib} в начале файла
\prog{bench/loop-select-bench.ps1}.

\section{Реализация функций конкатенации}
\label{app:concat-impl}

Четыре вспомогательных метода класса \prog{ExprOps}, на которые
опирается \prog{Expr.concat}. Каждый обрабатывает свою комбинацию
типов операндов; диспетчер показан в листинге~\ref{lst:concat-cases}
основного текста.

\begin{longlisting}
\caption{concatSS: оба операнда ShallowDeque}\label{lst:concat-ss}
\begin{minted}{java}
private static <T> Expr<T> concatSS(ShallowDeque<T> a, ShallowDeque<T> b) {
    int total = a.size() + b.size();
    if (total <= CAPACITY) {
        ShallowDeque<T> r = a;
        for (T x : iter(b)) {
            r = r.pushBackTight(x);
        }
        return r;
    }
    // total > CAPACITY: один из буферов полный или почти полный;
    // делаем новый Deep с пустой mid.
    return DeepDeque.mkDeep(a, ShallowDeque.empty(), b, total);
}
\end{minted}
\end{longlisting}

\begin{longlisting}
\caption{concatSD: ShallowDeque + DeepDeque}\label{lst:concat-sd}
\begin{minted}{java}
private static <T> Expr<T> concatSD(ShallowDeque<T> a, DeepDeque<T> b) {
    ShallowDeque<T> bf = b.front();
    int total = b.size() + a.size();
    // Сливаем 'a' слева в b.front. Если итог влезает в один Shallow,
    // он становится новым front; иначе расщепляем на два,
    // второй идёт в начало mid.
    int bfTotal = a.size() + bf.size();
    if (bfTotal <= CAPACITY) {
        ShallowDeque<T> mergedFront = ShallowDeque.empty();
        for (T x : iter(a))  mergedFront = mergedFront.pushBackTight(x);
        for (T x : iter(bf)) mergedFront = mergedFront.pushBackTight(x);
        return DeepDeque.mkDeep(mergedFront, b.mid(), b.back(), total);
    }
    // Иначе формируем плоский массив и делим примерно пополам.
    Object[] tmp = new Object[bfTotal];
    int i = 0;
    for (T x : iter(a))  tmp[i++] = x;
    for (T x : iter(bf)) tmp[i++] = x;
    int leftLen = bfTotal - CAPACITY;
    if (leftLen < ShallowDeque.MIN_FILL) leftLen = ShallowDeque.MIN_FILL;
    int rightLen = bfTotal - leftLen;
    ShallowDeque<T> head = ShallowDeque.ofArray(tmp, 0, leftLen);
    ShallowDeque<T> tail = ShallowDeque.ofArray(tmp, leftLen, rightLen);
    // head становится новым front; tail - первым бакетом mid.
    Expr<ShallowDeque<T>> newMid = b.mid().pushFront(tail);
    return DeepDeque.mkDeep(head, newMid, b.back(), total);
}
\end{minted}
\end{longlisting}

\begin{longlisting}
\caption{concatDS: DeepDeque + ShallowDeque}\label{lst:concat-ds}
\begin{minted}{java}
private static <T> Expr<T> concatDS(DeepDeque<T> a, ShallowDeque<T> b) {
    ShallowDeque<T> ab = a.back();
    int total = a.size() + b.size();
    int abTotal = ab.size() + b.size();
    if (abTotal <= CAPACITY) {
        ShallowDeque<T> mergedBack = ShallowDeque.empty();
        for (T x : iter(ab)) mergedBack = mergedBack.pushBackTight(x);
        for (T x : iter(b))  mergedBack = mergedBack.pushBackTight(x);
        return DeepDeque.mkDeep(a.front(), a.mid(), mergedBack, total);
    }
    Object[] tmp = new Object[abTotal];
    int i = 0;
    for (T x : iter(ab)) tmp[i++] = x;
    for (T x : iter(b))  tmp[i++] = x;
    int rightLen = abTotal - CAPACITY;
    if (rightLen < ShallowDeque.MIN_FILL) rightLen = ShallowDeque.MIN_FILL;
    int leftLen = abTotal - rightLen;
    ShallowDeque<T> head = ShallowDeque.ofArray(tmp, 0, leftLen);
    ShallowDeque<T> tail = ShallowDeque.ofArray(tmp, leftLen, rightLen);
    Expr<ShallowDeque<T>> newMid = a.mid().pushBack(head);
    return DeepDeque.mkDeep(a.front(), newMid, tail, total);
}
\end{minted}
\end{longlisting}

\begin{longlisting}
\caption{concatDD: два DeepDeque --- ключевой случай O(1) аморт.}\label{lst:concat-dd}
\begin{minted}{java}
private static <T> Expr<T> concatDD(DeepDeque<T> a, DeepDeque<T> b) {
    int total = a.size() + b.size();
    ShallowDeque<T> ab = a.back();
    ShallowDeque<T> bf = b.front();
    int boundary = ab.size() + bf.size();

    Expr<ShallowDeque<T>> aMidPlus = a.mid();
    if (boundary == 0) {
        // оба пограничных пусты - ничего не добавляем.
    } else if (boundary <= CAPACITY) {
        ShallowDeque<T> merged = ShallowDeque.empty();
        for (T x : iter(ab)) merged = merged.pushBackTight(x);
        for (T x : iter(bf)) merged = merged.pushBackTight(x);
        aMidPlus = aMidPlus.pushBack(merged);
    } else {
        Object[] tmp = new Object[boundary];
        int i = 0;
        for (T x : iter(ab)) tmp[i++] = x;
        for (T x : iter(bf)) tmp[i++] = x;
        int rightLen = boundary - CAPACITY;
        if (rightLen < ShallowDeque.MIN_FILL) rightLen = ShallowDeque.MIN_FILL;
        int leftLen = boundary - rightLen;
        ShallowDeque<T> m1 = ShallowDeque.ofArray(tmp, 0, leftLen);
        ShallowDeque<T> m2 = ShallowDeque.ofArray(tmp, leftLen, rightLen);
        aMidPlus = aMidPlus.pushBack(m1).pushBack(m2);
    }

    // Рекурсивное склеивание middle-очередей: O(1) аморт. по Окасаки.
    Expr<ShallowDeque<T>> newMid = concat(aMidPlus, b.mid());
    return DeepDeque.mkDeep(a.front(), newMid, b.back(), total);
}
\end{minted}
\end{longlisting}
